\documentclass[UTF8,a4paper,12pt,fontset=windows]{ctexart}

\usepackage[top=2.2cm,bottom=2.0cm,left=2.5cm,right=2.5cm]{geometry}
\usepackage{amsmath}
\usepackage{graphicx}
\usepackage{enumitem}
\usepackage{fancyhdr}
\usepackage{lastpage}
\usepackage{caption}

\setlength{\parindent}{0pt}
\setlength{\parskip}{0pt}
\linespread{1.35}
\allowdisplaybreaks

\pagestyle{fancy}
\fancyhf{}
\fancyfoot[C]{\small 第 \thepage\ 页，共 \pageref{LastPage} 页}
\renewcommand{\headrulewidth}{0pt}
\renewcommand{\footrulewidth}{0pt}

\newlist{examquestions}{enumerate}{1}
\setlist[examquestions]{
  label=\chinese{examquestionsi}、,
  labelsep=0.35em,
  leftmargin=2.6em,
  itemindent=0pt,
  labelwidth=2.2em,
  align=left,
  itemsep=1.15em,
  topsep=0pt,
  parsep=0pt
}

\begin{document}

\begin{center}
  {\zihao{4}\heiti 2026 年秋季学期}\\[0.65em]
  {\zihao{2}\heiti 数字集成电路设计期末试题}\\[0.45em]
  {\zihao{-4}回忆人：zsy}
\end{center}

\vspace{0.35em}
\rule{\textwidth}{0.8pt}
\vspace{1.2em}

\begin{examquestions}
  \item 以 CMOS 反相器为例，请写出 CMOS 集成制造工艺的步骤和每步的意义，并画出剖面图。

  \item 请阐述闩锁效应的原理，画出示意图，并说明如何抑制闩锁效应。

  \item 画出 CMOS 反相器的电压传输特性曲线，并说明曲线每一段对应的 PMOS、NMOS 状态以及 $V_i$ 的取值范围。写出逻辑转折电压 $V^{*}$ 的表达式并说明其意义。

  \item 画出静态 D 触发器的电路（逻辑图），说明保持/建立时间的作用。

  \item 一个 CMOS 反相器的 PMOS 的 $(W/L)_p=2/1$，NMOS 的 $(W/L)_n=1/1$，请你设计驱动能力不弱于该反相器的二输入与非门和或非门尺寸 $W/L$，比较三个电路输出高、低电平的噪声容限大小。

  \item 下图为同相施密特电路，请你描述其工作原理，并分析其抗噪声能力。

  \begin{center}
    \vspace{0.4em}
    \includegraphics[width=6.4cm]{assets/schmitt.png}
  \end{center}
\end{examquestions}

\end{document}
